\section{Comparison with Existing Frameworks}
\label{sec:compare}

%% tab-platforms.tex -- Table: feature-by-feature against the prompt-lifecycle
%% platforms, which are the genuinely closest systems.
%%
%% This table exists because an earlier draft asserted that the release concern
%% was unoccupied. It is not. MLflow's Prompt Registry, LangSmith, and Langfuse
%% all ship versioned prompts, a named pointer the client resolves at call time,
%% rollback by re-pointing, and role restrictions on who may move the pointer.
%% A comparison that gave those systems "none" for governance was a straw man,
%% and the delta this paper can honestly claim is narrower than the one the
%% earlier draft claimed.
%%
%% Every row is sourced to dated product documentation. Product capability moves;
%% the access date is part of the claim, and the caption says so.

\begin{table}[tbp]
\centering
\caption{Feature-by-feature comparison with the prompt-lifecycle platforms --- the
closest systems to \sys, and the ones its contribution should be judged against.
\textbf{Claims are as of 2026-08-03}, sourced to the cited documentation; these are
live products and a reader checking later should expect drift. Above the rule are
the four concerns an earlier draft of this paper got wrong: versioning, call-time
pointer resolution, rollback, and role-restricted promotion are \emph{present} in
all three, and \sys is not novel in any of them. Below the rule is where a
difference survives scrutiny. \secref{sec:compare:platforms} argues both halves.}
\label{tab:platforms}
\small
\setlength{\tabcolsep}{4.5pt}
\renewcommand{\arraystretch}{1.10}
\begin{cbwide}
\begin{tabular}{@{}L{2.86cm}L{3.16cm}L{3.16cm}L{3.16cm}VL{3.30cm}@{}}
\toprule
\theadrow
\thead{Concern} &
\thead{MLflow Prompt Registry} & \thead{LangSmith} & \thead{Langfuse} &
\thead{\sysbare{}} \\
\midrule

\textbf{Prompt versioning} &
immutable commit-based versions, with diffs~\cite{mlflowpromptregistry} &
commit history per prompt~\cite{langsmithprompts} &
version identifiers per prompt~\cite{langfuseprompts} &
\cellcolor{cbControlBg}versions per \gasset{} --- the same mechanism, not a
better one \\

\zebra
\textbf{Call-time pointer} &
mutable \emph{aliases}~\cite{mlflowpromptregistry} &
commit tags; reserved \code{production}, pulled by name at
runtime~\cite{langsmithprompts} &
\emph{labels}; an unlabelled fetch serves \code{production}~\cite{langfuseprompts} &
\cellcolor{cbControlBg}\livetarget{} per family --- \textbf{not a differentiator}
\\

\textbf{Rollback} &
re-point the alias~\cite{mlflowpromptregistry} &
rollback to a prior commit~\cite{langsmithprompts} &
re-label an older version~\cite{langfuseprompts} &
\cellcolor{cbControlBg}restore a \emph{recorded} checkpoint --- governed path only
(\secref{sec:limitations}) \\

\zebra
\textbf{Who may promote} &
registry permissions, where the deployment provides them &
owners-only tag create and delete~\cite{langsmithprompts} &
\emph{protected labels}, by project role~\cite{langfuseprompts} &
\cellcolor{cbControlBg}\operator{} or \admin{} scope --- comparable, and no
separation of duties (\secref{sec:limitations}) \\

\midrule

\textbf{Families beyond prompts} &
models; no other artifact type is a registry citizen &
datasets &
datasets &
\cellcolor{cbControlBg}\statFamiliesW{} families --- workflows, tools, MCP
servers, corpora, skills --- under one mode vocabulary
(\tabref{tab:families}) \\

\zebra
\textbf{Evidence bound to the version} &
evaluation integrated; the score annotates &
experiments and datasets, linked &
evaluations and scores, linked &
\cellcolor{cbControlBg}scores, corpus identity, diagnosis, and optimizer choice
bound to the candidate and carried into review \\

\textbf{Score as a precondition on promotion} &
none documented: promotion is a pointer move &
none documented &
none documented &
\cellcolor{cbControlBg}a failing gate stops the candidate --- inside the
refinement state machine, and only there \\

\zebra
\textbf{Automated candidate proposal} &
\textsc{gepa} and metaprompting
optimizers~\cite{mlflowoptimize} &
none documented &
none documented &
\cellcolor{cbControlBg}optimizers wired \emph{to} the gate --- the optimization is
not the delta, the wiring is \\

\textbf{Cross-family impact} &
\chipnone{} --- no second family to impact &
\chipnone{} &
\chipnone{} &
\cellcolor{cbControlBg}a change is relatable to dependents, because dependents are
records \\

\bottomrule
\end{tabular}
\end{cbwide}
\tabnote{``None documented'' means the mechanism was not found in the cited pages
on the access date --- a claim about the documentation, which a reader can check,
rather than about the product, which we are not positioned to make. Tracing-only
systems such as Phoenix~\cite{phoenix} are excluded: they do not offer prompt
lifecycle management and these rows would score them unfairly.}
\end{table}

%% tab-compare.tex -- Table: architectural comparison with adjacent systems.
%%
%% Paired with fig-capability. The figure answers "is the capability there?"; this
%% table answers "what is the architectural shape, and what does it cost?"
%%
%% Readability rules applied here:
%%   * \small, not \footnotesize. A comparison table nobody can read is not
%%     evidence, so the cells were shortened to pay for the larger type.
%%   * At most one clause per cell. The argument lives in Section 11; the table
%%     carries the shape, not the reasoning.
%%   * Alternating row tint, so the eye tracks a row across six columns.
%%   * The final column is not optional -- a comparison with no cost column is a
%%     marketing artifact.

\begin{table}[tbp]
\centering
\caption{Architectural comparison across the dimensions a systems reader asks
about. These systems are not substitutes, so the table compares \emph{shapes and
trade-offs} rather than ranking them; \secref{sec:compare} carries the reasoning,
and the last row states what each design gives up, including \sys's. This table
stops at shape: \tabref{tab:platforms} compares the prompt-lifecycle column
feature by feature against dated sources.}
\label{tab:compare}
\small
\setlength{\tabcolsep}{4pt}
\renewcommand{\arraystretch}{1.12}
\begin{cbwide}
\begin{tabular}{@{}L{2.20cm}VL{3.00cm}VL{2.72cm}L{2.68cm}L{2.52cm}L{3.08cm}@{}}
\toprule
\theadrow
\thead{Dimension} & \thead{\sysbare{}} &
\thead{Prompt lifecycle + obs.} & \thead{Composition} &
\thead{Multi-agent} & \thead{Optimizers + harnesses} \\
\theadrow
& &
{\footnotesize\itshape\color{white} \mlflow{}, LangSmith, Langfuse; Phoenix (tracing only)} &
{\footnotesize\itshape\color{white} LangChain, LangGraph, LlamaIndex} &
{\footnotesize\itshape\color{white} AutoGen, CrewAI} &
{\footnotesize\itshape\color{white} DSPy, promptfoo, Evals} \\
\midrule

\textbf{Organizing abstraction} &
\cellcolor{cbControlBg}the \gasset, across families &
the trace, plus the versioned prompt &
the chain or graph &
the agent and its conversation &
the program or test case \\

\zebra
\textbf{Architecture} &
\cellcolor{cbControlBg}layered control plane; one ASGI process &
service plus SDK &
library, inside the host app &
library or agent runtime &
compiler or CLI, out of band \\

\textbf{Modularity} &
\cellcolor{cbControlBg}narrow protocols at named seams &
exporters and hooks &
very high: composition \emph{is} the product &
agent, tool, memory interfaces &
optimizer and metric interfaces \\

\zebra
\textbf{Extensibility} &
\cellcolor{cbControlBg}new family, optimizer or scorer --- each must be
\emph{wired} &
scorers and dashboards &
components, freely &
agents and tools &
optimizers and metrics \\

\textbf{Memory \& execution} &
\cellcolor{cbControlBg}relational metadata authoritative; claim-and-lease queues &
append-mostly trace store &
in-process, host-provided &
conversation buffers &
batch, outside serving \\

\zebra
\textbf{Scalability shape} &
\cellcolor{cbControlBg}stateless replicas; the DB is the limit &
collector scales; the trace store dominates cost &
whatever the host is &
conversation length and token spend &
the optimizer's evaluation budget \\

\textbf{Developer experience} &
\cellcolor{cbControlBg}browser-first operate-and-release; big config surface &
excellent for diagnosis &
excellent for building &
excellent for prototyping &
excellent offline \\

\zebra
\textbf{Automation reach} &
\cellcolor{cbControlBg}trace\,\arrowto\,signal\,\arrowto\,candidate\,\arrowto\,gate\,\arrowto\,reviewed release &
alerting; \mlflow{} adds evaluate\,\arrowto\,optimize\,\arrowto\,alias, self-assembled &
none, by design &
task decomposition within a run &
candidate generation only \\

\textbf{Governance model} &
\cellcolor{cbControlBg}typed asset, live target, scopes, audit --- per family &
for prompts: version, alias, role-restricted promotion --- but the score does not
gate~\cite{mlflowpromptregistry,langfuseprompts} &
none &
ad~hoc per conversation &
none \\

\midrule
\textbf{What it gives up} &
\cellcolor{cbControlBg}\textbf{not} a composition or multi-agent framework; one
environment; no managed cloud; per-family guarantees &
governs prompts and models only; no score gates the promotion &
no version, evidence or release semantics; every team rebuilds them &
no governed release path; provenance is ad~hoc &
no live target, no audit trail, no human decision point \\

\bottomrule
\end{tabular}
\end{cbwide}
\tabnote{Rules set off the \sysbare{} column; the tints and bands are reading aids
and carry no other meaning.}
\end{table}


This section makes two comparisons, and separating them matters. The first is
narrow and is where the paper is most exposed: against the platforms that already
manage a prompt lifecycle. The second is broad and is about architectural shape.
Doing only the second would let a shape argument stand in for a feature argument,
which is how a comparison section becomes advertising.

\subsection{The prompt-lifecycle platforms already do much of this}
\label{sec:compare:platforms}

An earlier version of this paper claimed the release concern was unoccupied. That
was wrong, and \tabref{tab:platforms} states it plainly. As of 2026-08-03:
\mlflow{}'s Prompt Registry provides commit-based, immutable prompt versions with
diffs and mutable \emph{aliases}~\cite{mlflowpromptregistry}; LangSmith provides
commit history, custom commit tags, reserved \code{staging} and \code{production}
tags resolved at pull time, documented rollback to a prior commit, and owners-only
control over tagging~\cite{langsmithprompts}; and Langfuse provides versions,
\code{production} and custom \emph{labels} resolved client-side at fetch time,
rollback by re-assigning a label, and \emph{protected labels} restricted by
role~\cite{langfuseprompts}. \mlflow{} additionally ships prompt optimization ---
\code{optimize\_prompts()} over \textsc{gepa} and metaprompting, driven by a
labelled dataset and scorers~\cite{mlflowoptimize}.

Four claims a reader might expect this paper to make are therefore unavailable to
it. Late-bound pointer indirection for prompts is not novel; neither is rollback by
re-pointing; neither is restricting who may promote; neither is generating a
candidate automatically. Where an earlier draft treated these as \sys's
contribution, they are table stakes, and the honest reading of the top block of
\tabref{tab:platforms} is that \sys re-implements a solved problem for prompts.

\para{What survives.} Three differences do survive, and they are narrower than the
old framing but they are checkable. \emph{First, family heterogeneity.} These
platforms govern prompts, and \mlflow{} also governs models; workflows, tool
definitions, MCP server configurations, and retrieval corpora are not registry
citizens with versions and pointers. \sys's contribution is the claim that one
mode vocabulary can span \statFamiliesW{} families --- together with the finding,
which is the paper's actual intellectual content, that spanning them costs uniform
guarantees (\secref{sec:arch:claim}). \emph{Second, evidence as a precondition
rather than an annotation.} All three platforms link evaluations to prompts; none
of their documentation describes a mechanism that makes a promotion \emph{fail}
because a score fell. In \sys the advancement gate is such a mechanism --- inside
the refinement state machine, and, as \secref{sec:limitations} concedes, only
there. \emph{Third, the wiring rather than the parts.} \mlflow{} optimizes
prompts and \mlflow{} versions prompts, but the path from a flagged production
trace to an optimizer run to a gated candidate to a reviewed release is assembled
by the adopter. \sys's claim is that this path is worth being a first-class
architectural object.

That is a smaller claim than ``the concern is unoccupied,'' and we would rather
state the smaller one accurately. A reviewer who finds even this delta insufficient
is making a reasonable judgement about significance rather than catching an error.

\subsection{Why a ranking across the broader field would be the wrong output}

The four groups in \tabref{tab:compare} are not substitutes. A team does not choose
between LangGraph and \sys any more than it chooses between a web framework and a
deployment pipeline; the realistic configuration is an agent composed in a framework,
traced by an observability platform, optimized by a program optimizer, and governed
by a control plane. Producing a single ranking across them would require a common
objective they do not share.

What can be compared is the \emph{organizing abstraction}, because that choice
determines what each system can and cannot express. Tracking and observability
platforms organize around the trace \emph{and}, in the three cases discussed
above, around a versioned prompt --- so they are authoritative about what happened
and, for prompts specifically, about what is live. What their organizing
abstraction does not reach is a second artifact type: a system whose registry
citizen is the prompt has no place to put a workflow definition. Composition
frameworks organize around the chain or graph, which makes them excellent for
construction and structurally unable to carry release semantics --- the graph is
constructed in the host process, so there is no artifact to version. Multi-agent
frameworks organize around the agent and the conversation, which makes provenance
per-conversation rather than per-artifact. Optimizers organize around the program
or the test case, which yields candidates and leaves the release to the caller.
\sys organizes around the \gasset, which is what lets one vocabulary carry
version, evidence, authority, and a live target across families --- and is also
why it has nothing useful to say about how to compose an agent.

\subsection{The dimensions where \sys is genuinely different}

\para{Memory and execution model.} This is the sharpest architectural difference.
The observability systems are append-mostly stores queried after the fact; the
composition and multi-agent frameworks hold state in process for the duration of a
run; the optimizers are batch processes outside the serving path. \sys is the only
system in the table for which \emph{relational metadata is authoritative} and for
which durable queue arbitration is part of the architecture rather than an
application concern. That is a consequence of being a control plane: a system that
must answer ``who approved this and when'' cannot hold that state in a process.

\para{Automation reach.} Composition frameworks reach, by design, nowhere --- they
are libraries. Multi-agent frameworks reach into task decomposition inside a single
run. Standalone optimizers reach to candidate generation and stop. The
prompt-lifecycle platforms reach further than the previous paragraph's framing
suggests: \mlflow{} in particular spans evaluation, optimization, versioning, and
alias promotion, so its reach ends at an \emph{operator-assembled} release rather
than at alerting~\cite{mlflowpromptregistry,mlflowoptimize}. \sys reaches from a
trace to a verified signal to a candidate to a scored gate to a reviewed release.
The distinguishing property is not the length of that span but that it is a single
wired path with a gate in it, rather than a set of capabilities an adopter connects
--- and \secref{sec:compare:platforms} is where that distinction is argued rather
than asserted.

\para{Extensibility, with a caveat we owe the reader.} The composition frameworks are
more extensible than \sys in the ordinary sense: adding a component is trivial and
requires no coordination. \sys's extensibility is narrower and has a specific
qualification --- adding a family requires \emph{explicitly wiring} its modes and
governance, per \secref{sec:arch:claim}. A reader comparing extensibility columns
should read \sys's as ``extensible at named seams, with an obligation attached,''
not as ``more extensible.''

\subsection{The dimensions where \sys is behind}

\para{Developer experience for construction.} If the task is to build an agent, a
composition framework is the better tool and \sys is not a substitute. \sys's
developer experience is browser-first and oriented toward the operate-and-release
loop; its configuration surface is large, which is a genuine onboarding cost that a
library does not impose.

\para{Multi-agent orchestration.} AutoGen and CrewAI are purpose-built for
composing agents into a running system, and \sys is not competitive there
(\figref{fig:capability}). Automatic optimization of a multi-agent \emph{bundle} ---
as opposed to a single prompt or skill --- is not a shipped \sys capability.

\para{Managed operation.} Every observability platform in the table offers a hosted
option. \sys must be run by the adopting team, on their own PostgreSQL, object
storage, and \mlflow{}. For a small team this may outweigh every architectural
consideration in this paper, and a comparison that omitted it would be dishonest.

\para{Uniformity.} An adopter who wants one release contract across all artifact
types will find \sys's per-family guarantees unsatisfying. \secref{sec:arch:claim}
argues this is the honest factoring rather than an incompleteness, but the argument
does not eliminate the cost: understanding \sys requires reading
\tabref{tab:families} rather than one sentence, and that is a real burden on the
adopter.

\subsection{Positioning, stated plainly}

\sys's claim is not that it does more than these systems, and it is not that the
release path is unoccupied --- for prompts, it is occupied, by at least three
products, with mechanisms that closely resemble \sys's
(\secref{sec:compare:platforms}). The claim is narrower on two axes. First, that
the release path generalizes \emph{beyond prompts} to workflows, tools, MCP
servers, and corpora, and that a single mode vocabulary can span them. Second, that
what the generalization costs --- guarantees that are per-family rather than
uniform --- is a result worth reporting rather than a defect to be hidden.

Whether this warrants a separate control plane, or should be absorbed into a
tracking platform, is a legitimate open question, and the evidence currently tilts
toward absorption being feasible: \mlflow{} already has the registry, the aliases,
the evaluation harness, and the optimizer, and what it lacks is a second artifact
type and a gate. An absorbing system would need to add a typed multi-family
registry, an authority model that distinguishes authoring from promoting, and an
enforced precondition on the pointer move. That is a smaller distance than an
earlier draft of this paper implied.
